\chapterbackmatter{Supplementary Exercises}

\begin{enumerate}
  \SupplementaryExercise{1}{Graded Algebra and Supersymmetry Transformations}
    {(Inspired by Haag, {\L}opusza\'nski, and Sohnius, ``All Possible Generators of Supersymmetries of the S-Matrix'', Sections 2--3; independently formulated.)}
    {ch25-s01-graded-algebra-and-supersymmetry-transformations}
  \begin{enumerate}
    \item \emph{The Supertrace Test.}
    Let \(V=V_{\bar0}\oplus V_{\bar1}\), and write endomorphisms in even and
      odd block form.  Define
      \(\operatorname{str}\begin{psmallmatrix}A&B\\C&D\end{psmallmatrix}
      =\operatorname{tr}A-\operatorname{tr}D\).
      Prove for homogeneous \(X,Y\) that
      \(\operatorname{str}[X,Y\}=0\), and conclude that the supertraceless
      matrices close under the graded bracket.
    \item \emph{The Signed Jacobi Identity.}
    For homogeneous elements of an associative graded algebra, set
      \([X,Y\}=XY-(-1)^{|X||Y|}YX\).  Expand all products and verify
      \[
        (-1)^{|X||Z|}[X,[Y,Z\}\}
        +(-1)^{|Y||X|}[Y,[Z,X\}\}
        +(-1)^{|Z||Y|}[Z,[X,Y\}\}=0.
      \]
      Explain why the case of three odd elements involves three
      commutators with anticommutators.
    \item \emph{Why a SUSY Variation Is Even.}
    Let \(Q_\alpha\) and a constant parameter \(\epsilon^\alpha\) both be
      odd.  Show that \(G=\epsilon^\alpha Q_\alpha+
      \epsilon^\dagger_{\dot\alpha}\bar Q^{\dot\alpha}\) is even and that
      \(\delta_\epsilon X=i[G,X]\) preserves the statistics of a homogeneous
      operator \(X\).  Then explain why treating \(\epsilon\) as commuting
      gives the wrong bracket in the commutator of two transformations.
  \end{enumerate}

  \SupplementaryExercise{2}{Spinorial Lorentz Generators}
    {(Adapted from David Tong, \emph{Supersymmetry Example Sheet 1}, Question 4.)}
    {ch25-s02-spinorial-odd-generators}
Use the two-component matrices implicit in Eq.~\eqref{eq:25.2.37},
  \[
    s^\mu=(\mathbf{1},\boldsymbol{\sigma}),\qquad
    \bar s^\mu=(\mathbf{1},-\boldsymbol{\sigma}),
    \qquad
    s^\mu\bar s^\nu+s^\nu\bar s^\mu
      =-2\eta^{\mu\nu}\mathbf{1} .
  \]
  In the mostly-plus convention of this volume, define
  \[
    \Sigma^{\mu\nu}
      =\frac{i}{4}\left(s^\mu\bar s^\nu
                       -s^\nu\bar s^\mu\right).
  \]
  Then use the two-component identities of Section 25.2 to prove directly that
  these matrices furnish the left-handed spinor representation of the
  Lorentz algebra:
  \[
    [\Sigma^{\mu\nu},\Sigma^{\rho\sigma}]
      =-i\left(
        \eta^{\nu\rho}\Sigma^{\mu\sigma}
        -\eta^{\nu\sigma}\Sigma^{\mu\rho}
        +\eta^{\mu\sigma}\Sigma^{\nu\rho}
        -\eta^{\mu\rho}\Sigma^{\nu\sigma}
      \right).
  \]
  In particular, derive rather than assume the useful intermediate identity
  \(\Sigma^{\mu\nu}
    =-\tfrac{i}{2}(\eta^{\mu\nu}\mathbf{1}
                   +s^\nu\bar s^\mu)\).

  \SupplementaryExercise{3}{The Superspin Casimir}
    {(Adapted from University of Cambridge, 2025 Part III Paper 307, \emph{Supersymmetry}, Question 3.)}
    {cambridge-part-iii-supersymmetry-2025-p307-q3-v3ch25-p03}
The generators of the Poincar\'e group \(M^{\mu\nu}\) and \(P^\sigma\)
  satisfy
  \[
    [M^{\mu\nu},P^\sigma]
      =i\bigl(P^\mu\eta^{\nu\sigma}-P^\nu\eta^{\mu\sigma}\bigr).
  \]
  Define the Pauli--Lubanski operator \(W_\mu\) in terms of
  \(M^{\mu\nu}\) and \(P^\mu\).

  Write down a general Poincar\'e transformation on a left-handed spinor
  field \(\psi_a(x)\) in terms of the left-handed spinor representation of
  the Lorentz algebra
  \[
    (\Sigma^{\mu\nu})^b{}_a
      =\frac{i}{4}
        \bigl(s^\mu\bar s^\nu-s^\nu\bar s^\mu\bigr)^b{}_a .
  \]
  Here
  \(s^\mu=(\mathbf 1,\boldsymbol{\sigma})\) and
  \(\bar s^\mu=(\mathbf 1,-\boldsymbol{\sigma})\).  In the mostly-plus
  convention of this volume, the self-duality relation supplied in the
  source becomes
  \[
    \Sigma^{\mu\nu}
      =\frac{1}{2i}\epsilon^{\mu\nu\rho\sigma}\Sigma_{\rho\sigma}.
  \]

  In terms of the
  \((SU(2)_{\mathrm L},SU(2)_{\mathrm R})\) representation spaces of the
  Lorentz algebra, \(\psi\sim(\tfrac12,0)\).  Write down
  \((\tfrac12,0)\otimes(\tfrac12,0)\) in terms of irreducible
  \((SU(2)_{\mathrm L},SU(2)_{\mathrm R})\) representation spaces.  Show
  this explicitly by decomposing \(\psi_a\chi_b\), where
  \(\chi\sim(\tfrac12,0)\).

  Write down all commutators or anticommutators of the super-Poincar\'e
  algebra that involve \(Q_a\) and/or \(\overline Q_{\dot a}\).  Then do
  the following.
  \begin{enumerate}
    \item Defining
      \[
        B_\mu=W_\mu-\frac14
          \bigl(\overline Q\,\bar s_\mu Q\bigr),
      \]
      calculate \([B_\mu,P_\rho]\).
    \item Defining
      \[
        C_{\mu\nu}=B_\mu P_\nu-B_\nu P_\mu,
      \]
      calculate \([C_{\mu\nu},P_\rho]\).
    \item Compute
      \(\bigl[\overline Q\,\bar s_\mu Q,Q_a\bigr]\).
    \item Compute \([B_\mu,Q_a]\).
    \item Compute \([C_{\mu\nu},Q_a]\).
    \item Defining the Lorentz-scalar superspin operator
      \[
        \widetilde C_2=C_{\mu\nu}C^{\mu\nu},
      \]
      write down the implied commutator with \(M^{\mu\nu}\).  Thus
      demonstrate that \(\widetilde C_2\) is a Casimir of the
      super-Poincar\'e algebra.
  \end{enumerate}

  \SupplementaryExercise{4}{Antisymmetric Central Charges and Their Normal Form}
    {(Inspired by Adel Bilal, \emph{Introduction to Supersymmetry}, Sections 3.1--3.2 and 8; independently formulated.)}
    {ch25-s04-antisymmetric-central-charges-and-their-normal}
  \begin{enumerate}
    \item \emph{Antisymmetry of \(Z_{rs}\).}
    In the extended algebra write
      \(\{Q_{\alpha r},Q_{\beta s}\}
      =\epsilon_{\alpha\beta}Z_{rs}\).
      Exchange the two complete labels \((\alpha,r)\) and \((\beta,s)\) and
      prove \(Z_{rs}=-Z_{sr}\).  Deduce that simple supersymmetry admits no
      scalar central charge of this form.
    \item \emph{Canonical Blocks of the Central Charge.}
    A complex antisymmetric \(N\times N\) matrix can be brought by a unitary
      congruence to \(2\times2\) blocks
      \(z_i\begin{psmallmatrix}0&1\\-1&0\end{psmallmatrix}\), plus a zero block
      when \(N\) is odd.  Use this form to find the eigenvalues of
      \(Z^\dagger Z\), their multiplicities, and
      \(\operatorname{Tr}\sqrt{Z^\dagger Z}\).
  \end{enumerate}

  \SupplementaryExercise{5}{The \(N=1\) Algebra, Parity, and State Pairing}
    {(Adapted from University of Cambridge, 2024 Part III Paper 307, \emph{Supersymmetry}, Question 1.)}
    {cambridge-part-iii-supersymmetry-2024-p307-q1-v3ch25-p05}
The generators of the Poincar\'e algebra are the Hermitian operators
  \(M^{\mu\nu}\) and \(P^\sigma\), with algebra
  \[
  \begin{aligned}
    [P^\mu,P^\nu]&=0,\\
    [M^{\mu\nu},P^\sigma]
      &=i\bigl(P^\mu\eta^{\nu\sigma}
               -P^\nu\eta^{\mu\sigma}\bigr),\\
    [M^{\mu\nu},M^{\rho\sigma}]
      &=i\bigl(
          M^{\mu\sigma}\eta^{\nu\rho}
         +M^{\nu\rho}\eta^{\mu\sigma}
         -M^{\mu\rho}\eta^{\nu\sigma}
         -M^{\nu\sigma}\eta^{\mu\rho}
        \bigr).
  \end{aligned}
  \]
  A spinor representation of the Lorentz algebra is provided by
  \[
    (\Sigma^{\mu\nu})^b{}_a
      =\frac{i}{4}
        \bigl(s^\mu\bar s^\nu-s^\nu\bar s^\mu\bigr)^b{}_a,
  \]
  where
  \[
    (s^\mu)_{a\dot a}
       =(\mathbf 1,\sigma^1,\sigma^2,\sigma^3)_{a\dot a},
    \qquad
    (\bar s^\mu)^{\dot a a}
       =(\mathbf 1,-\sigma^1,-\sigma^2,-\sigma^3)^{\dot a a}.
  \]
  The displayed relations are written in the mostly-plus metric of this
  volume.
  \begin{enumerate}
    \item In one sentence, describe what the Coleman--Mandula theorem says
      about extensions of the Poincar\'e algebra in quantum field theory.
      Briefly discuss the assumption that can be weakened to allow
      supersymmetry, including a short description of the type of algebra
      that results.
    \item What does the Coleman--Mandula theorem imply for the algebra
      between an internal-symmetry generator \(T_A\) and a supersymmetry
      generator \(Q_a\)?  Ignoring \(R\)-symmetries, extend the
      Poincar\'e algebra with the \(N=1\) generators
      \(Q_a,\overline Q_{\dot a}\), giving detailed arguments for the
      chosen form of each relation.
    \item Consider the action under a parity-reversal operator
      \(\widehat{\mathsf P}\), where \(|\eta_{\mathsf P}|=1\) and
      \[
      \begin{aligned}
        \widehat{\mathsf P}Q_a\widehat{\mathsf P}^{-1}
          &=\eta_{\mathsf P}(s^0)_{a\dot b}
            \overline Q^{\dot b},\\
        \widehat{\mathsf P}\overline Q^{\dot a}
          \widehat{\mathsf P}^{-1}
          &=-\eta_{\mathsf P}^{*}
            (\bar s^0)^{\dot a b}Q_b .
      \end{aligned}
      \]
      By performing a parity transformation on each side of the algebra
      between \(Q_a\) and \(\overline Q_{\dot b}\) that you derived in
      part (b), check that the relation is consistent with a parity
      transformation.
    \item Define the fermion-number operator \((-1)^F\) by its action on
      bosonic states \(\lvert B\rangle\) and fermionic states
      \(\lvert F\rangle\).  Prove that the number of fermionic states
      \(n_F\) equals the number of bosonic states \(n_B\) in any
      supermultiplet.  Where does your argument break down for theories
      that incorporate additional supersymmetry-breaking terms in the
      Lagrangian?
  \end{enumerate}

  \SupplementaryExercise{6}{Massive \(N=2\) Multiplets and the BPS Bound}
    {(Adapted from University of Cambridge, 2017 Part III Paper 307, \emph{Supersymmetry}, Question 1.)}
    {cambridge-susy-2017-q1-v3ch25-s06}
  \begin{enumerate}
    \item Write the four-dimensional extended supersymmetry algebra for
      \(\WeylQ_{ar}\), \(r=1,2\), including its antisymmetric central charge
      \(Z_{rs}\), the Lorentz brackets, and the commutators with momentum.
      Work in the normalization of Eqs.~\eqref{eq:25.2.7}--\eqref{eq:25.2.10}.
    \item Set \(Z_{rs}=0\), go to the rest frame of a massive representation,
      and normalize the supercharges as fermionic creation and annihilation
      operators.  Starting from a spin-zero Clifford vacuum, construct the
      full representation.  Give the spin and multiplicity at every fermion
      number and verify the total number of states.
    \item Now take \(Z_{12}=z\ne0\).  Rephase the supercharges so that \(z\)
      is real and nonnegative, and form linear combinations of
      \(\WeylQ_{a1}\) and \(e_{ab}\WeylQ_{b2}^\dagger\) whose positive
      anticommutators are proportional to \(2M+|z|\) and \(2M-|z|\).
      Derive the BPS inequality in this volume's convention.  At saturation,
      identify the null operators, construct the shortened representation
      based on a spin-zero Clifford vacuum, and state the fraction of
      supersymmetry preserved.
    \item Generalize the construction to even \(N\).  Put the complex
      antisymmetric central-charge matrix into unitary skew-block normal form,
      derive the bound associated with every skew eigenvalue, and determine
      the dimension of a massive representation when exactly \(k\) skew
      blocks saturate their bounds.
  \end{enumerate}

  \SupplementaryExercise{7}{Massive \(N=1\) Multiplets and Parity}
    {(Inspired by Adel Bilal, \emph{Introduction to Supersymmetry}, Sections 3.2--3.4; independently formulated.)}
    {ch25-s07-massive-multiplets-and-parity}
  \begin{enumerate}
    \item \emph{A Long Massive \(N=1\) Multiplet.}
    Start from a spin-\(j\) Clifford vacuum annihilated by the two rest-frame
      annihilation operators.  Apply the two creation operators and decompose
      the one-creation level under rotations.  For \(j>0\), determine every spin
      and its multiplicity, and verify equality of bosonic and fermionic
      degrees of freedom for both integral and half-integral \(j\).
    \item \emph{The Collapsed Massive Multiplet.}
    Repeat the massive \(N=1\) construction for a spin-zero Clifford vacuum.
      Explain why the formal \(j-1/2\) representation is absent, and show that
      the spectrum consists of two spin-zero particles and one spin-\(1/2\)
      particle.  Count physical spin states.
    \item \emph{Parity Along the Massive Ladder.}
    Assume parity is conserved and choose the phase convention
      \(P^{-1}Q_aP=i\epsilon_{ab}Q_b^\dagger\).  If the one-creation states of
      spin \(j\pm1/2\) have parity \(\eta\), prove that the two spin-\(j\)
      combinations have parities \(+i\eta\) and \(-i\eta\).  Specialize the
      result to the collapsed \(j=0\) multiplet.
  \end{enumerate}

  \SupplementaryExercise{8}{Massless Oscillators, Helicity, and CPT}
    {(Inspired by David Tong, \emph{Supersymmetric Field Theory}, Sections 2.3.2 and 2.4.1; independently formulated.)}
    {ch25-s08-massless-oscillators-helicity-and-cpt}
  \begin{enumerate}
    \item \emph{The Massless-Frame Oscillator.}
    Take \(p^\mu=(E,0,0,E)\) in a convention where
      \(\{Q_\alpha,Q^\dagger_{\dot\beta}\}
      =2\sigma^\mu_{\alpha\dot\beta}p_\mu\).  Diagonalize the resulting
      \(2\times2\) matrix.  Show that one component of \(Q\) has zero norm and
      hence acts trivially, while the other becomes a single fermionic
      oscillator.  State how its adjoint changes helicity.
    \item \emph{The Exterior Algebra of Helicity.}
    For \(N\) massless supercharges, identify the state space generated from a
      highest-helicity state with \(\bigwedge^\bullet\mathbb C^N\).
      Derive the multiplicity \(\binom Nk\) and helicity
      \(\lambda_{\max}-k/2\) at level \(k\).  Prove that the total number of
      states is \(2^N\).
    \item \emph{The CPT Self-Conjugacy Test.}
    A massless multiplet runs from \(\lambda_{\max}\) to
      \(\lambda_{\max}-N/2\).  Derive the condition for this helicity interval
      to be invariant under CPT.  If it is not invariant, write the top
      helicity of the required CPT-conjugate multiplet.
  \end{enumerate}

  \SupplementaryExercise{9}{The \(N=1\) Chiral and Vector Multiplets}
    {(Inspired by Stephen P. Martin, \emph{A Supersymmetry Primer}, version 7, Sections 1, 3.1, and 3.3; independently formulated.)}
    {ch25-s09-the-chiral-and-vector-multiplets}
  \begin{enumerate}
    \item \emph{The Massless Chiral Multiplet.}
    Construct the CPT-complete massless \(N=1\) multiplet containing a
      left-handed Weyl fermion.  List the helicities before and after CPT
      completion, interpret the result as one complex scalar plus one Weyl
      fermion, and count on-shell bosonic and fermionic degrees of freedom.
    \item \emph{The Massless Vector Multiplet.}
    Starting from helicity \(+1\), construct the \(N=1\) massless multiplet
      and its CPT conjugate.  Show that the completed spectrum is a massless
      vector plus a Majorana gaugino, and verify the equality of physical
      bosonic and fermionic polarizations.
  \end{enumerate}

  \SupplementaryExercise{10}{Extended Multiplets and the \(N=4\) \(N=1\) Decomposition}
    {(Adapted from John McGreevy, \emph{MIT 8.821 Problem Set 2}, Problem 4.)}
    {ch25-s10-extended-multiplets-and-the-decomposition}
  \begin{enumerate}
    \item \emph{Field Content in \(N=1\) Language.}
    Choose one \(N=1\) subalgebra of four-dimensional \(N=4\) super
      Yang--Mills.  Decompose the complete \(N=4\) vector multiplet into
      irreducible \(N=1\) supermultiplets, specifying the gauge representation
      and the component fields of each.
    \item \emph{The \(N=1\) Superspace Action.}
    Write the gauge kinetic term, the gauge-covariant Kähler term, and the
      superpotential in \(N=1\) superspace.  Fix the relative cubic coupling
      required by \(N=4\) supersymmetry and show that eliminating the auxiliary
      fields produces the commutator-squared scalar potential.
    \item \emph{The Hidden \(SO(6)\) Symmetry.}
    Reorganize the three complex adjoint scalars into six real scalars and the
      four Weyl fermions into a \(\mathbf4\) of
      \(SU(4)\simeq\operatorname{Spin}(6)\).
      Explain how the component interactions display this \(SO(6)\)
      \(R\)-symmetry even though only \(SU(3)\times U(1)\) is manifest in the
      chosen superspace description.  Finally, decide whether exact \(N=1\)
      supersymmetry together with this full \(R\)-symmetry is sufficient to
      recover all four supersymmetries.
  \end{enumerate}

  \SupplementaryExercise{11}{The \(N=8\) Gravity Multiplet and Helicity Cutoffs}
    {(Inspired by David Tong, \emph{Supersymmetric Field Theory}, Sections 2.2.1 and 2.3--2.4; independently formulated.)}
    {ch25-s11-the-gravity-multiplet-and-helicity-cutoffs}
  \begin{enumerate}
    \item \emph{The \(N=8\) Gravity Multiplet.}
    Starting from helicity \(+2\), use the exterior-algebra construction to
      list the complete \(N=8\) massless supergravity multiplet.  Identify the
      spin-2, spin-\(3/2\), spin-1, spin-\(1/2\), and spin-0 multiplicities and
      verify CPT self-conjugacy.
    \item \emph{Helicity Cutoffs on Extended SUSY.}
    Use the helicity width \(N/2\) of a massless multiplet to prove that a
      CPT-self-conjugate theory with no particles of spin above \(1\) has
      \(N\leq4\), while one with no particles of spin above \(2\) has
      \(N\leq8\).  State the assumptions under which these are bounds on global
      supersymmetry and supergravity, respectively.
  \end{enumerate}

  \SupplementaryExercise{12}{Protected Helicity Sums and Central-Charge Parity}
    {(Inspired by David Tong, \emph{Supersymmetric Field Theory}, Sections 2.3--2.4; independently formulated.)}
    {ch25-s12-protected-helicity-sums-and-central-charge}
  \begin{enumerate}
    \item \emph{Helicity Supertraces.}
    For a massless \(N\)-extended multiplet define
      \(B_m=\sum_{\text{states}}(-1)^F\lambda^m\).
      Using the generating function
      \(\sum(-1)^F e^{t\lambda}
      =e^{t\lambda_{\max}}(1-e^{-t/2})^N\), prove that
      \(B_m=0\) for \(m<N\).  Compute \(B_N\).
    \item \emph{Parity and the Central-Charge Phase.}
    With
      \(P^{-1}Q_{\alpha r}P=i\epsilon_{\alpha\beta}
      Q^\dagger_{\dot\beta r}\) and
      \(\{Q_{\alpha r},Q_{\beta s}\}
      =\epsilon_{\alpha\beta}Z_{rs}\), determine the parity transform of
      \(Z_{rs}\).  Explain when a single BPS central-charge sector can carry
      parity eigenstates and when parity must exchange two sectors.  Relate this
      to the freedom to rephase the supercharges.
  \end{enumerate}

\end{enumerate}
